% ==============================================================================
% Instituto Federal de Alagoas
% Modelo usado para Trabalho De Conclusão de Curso (TCC) da Coordenação de 
% Informática do IFAL - Campus Maceió
%
% Autor:  Leonardo Melo de Medeiros
% Endereço: https://www.overleaf.com/read/swrgmkkmhwzs
% Baseado no modelo do IFRN feito pelo Prof. Rodrigo Siqueira
% ==============================================================================

% Classe .......................................................................
\documentclass[a4paper,12pt,openany,oneside]{book}

%ATENÇÃO USAR PADRÃO DE CONDIFICAÇÃO DO LINUX UTF8
%PARA SABER SE VC O ESTÁ USANDO BASTA PERCEBER A ACENTUAÇÃO ERRADA
\usepackage[utf8]{inputenc}

%Para adicionar figuras
\usepackage{graphicx}



% Pacotes ....................................................................................
% Pacotes Principais -----------------------------------------------------------
\usepackage[portuges,brazil]{babel}
\usepackage{ae} 

%Linux ou Windows com UTF8
\usepackage[utf8]{inputenc}

%\usepackage[latin1]{inputenc}

%NO MAC ao invŽs de latin1 deve-se usar o applemac
%\usepackage[applemac]{inputenc}


% Formatacaoo de capitulos ------------------------------------------------------
%\usepackage{capitulos}

% Figuras e Imagens ------------------------------------------------------------
\usepackage{graphicx}
% Figuras lado a lado
%\usepackage{epsfig} - REMOVIDO
%\usepackage{subfigure} - REMOVIDO

% Para poder ter tabelas com colunas de largura auto-ajust‡vel
\usepackage{tabularx}

% Utilizar H para inserir as imagens REALMENTE onde eu desejo
%\usepackage{float} - REMOVIDO

% Para que os primeiros parágrafos das seções também sejam indentados
% \usepackage{indentfirst} - REMOVIDO

%Permitir usar o comando \citeasnoun que coloca o nome na cita?‹o
\usepackage{harvard}


% Fontes -----------------------------------------------------------------------
%\usepackage[T1]{fontenc} - REMOVIDO
%\usepackage{pslatex} - REMOVIDO

% Matemático -------------------------------------------------------------------
%\usepackage{amsmath} - REMOVIDO
%\usepackage{amstext} - REMOVIDO

% Simbolos ---------------------------------------------------------------------
%\usepackage{textcomp} - REMOVIDO

% Tabelas ----------------------------------------------------------------------
%\usepackage{multirow} - REMOVIDO
% Colorir a tabela
%\usepackage{colortbl} - REMOVIDO

% Glossário --------------------------------------------------------------------
%\usepackage[portuguese,noprefix]{nomencl}
%\usepackage{makeglo}

% Outros pacotes ---------------------------------------------------------------
%\usepackage{noitemsep} - REMOVIDO

%Subsubsection
%\usepackage{chngcntr} - REMOVIDO

% Para permitir espaçamento simples, 1 1/2 e duplo
\usepackage{setspace}

% Para executar um comando depois do fim da página corrente
%\usepackage{afterpage} - REMOVIDO

% Para poder incluir arquivos Postscript com cores (do Xfig, por exemplo)
%\usepackage{color} - REMOVIDO

  % Itens numerados
%\usepackage{enumerate} - REMOVIDO

% Comentários em bloco
%\usepackage{verbatim} - REMOVIDO

% Referências ------------------------------------------------------------------
%\usepackage{html} - REMOVIDO
%\usepackage{url} - REMOVIDO
%\usepackage[abbr]{harvard}	% As chamadas são sempre abreviadas - REMOVIDO
%\harvardparenthesis{square}	% Colchetes nas chamadas - REMOVIDO
%\harvardyearparenthesis{round}	% Parêntesis nos anos das referências - REMOVIDO
%\renewcommand{\harvardand}{e}	% Substituir "&" por "e" nas referências
%\renewcommand{\harvardurl}{URL: \url} - REMOVIDO

% Novos Comandos .....................................................................
\input{configs/comandos.tex}
%\citationmode{abbr}

% Inicia o texto
\begin{document}

% Paginas iniciais (sem numeração)
\pagestyle{empty}

% P?gina de rosto (capa interna)
%
% ********** Página de Rosto
%

% titlepage gera páginas sem numeração
\begin{titlepage}

\begin{center}

\small

% O comando @{} no ambiente tabular x é para criar um novo delimitador
% entre colunas que não a barra vertical | que é normalmente utilizada.
% O delimitador desejado vai entre as chaves. No exemplo, não há nada,
% de modo que o delimitador é vazio. Este recurso está sendo usado para
% eliminar o espaço que geralmente existe entre as colunas
\begin{tabularx}{\linewidth}{ c X }
% A figura foi colocada dentro de um parbox para que fique verticalmente
% centralizada em relação ao resto da linha
\parbox[c]{3cm}{\includegraphics[width=\linewidth]{IFRN}} &
\begin{center}
\textsf{\textsc{Instituto Federal de Alagoas\\
 Campus Maceió\\
 Coordenação de Informática \\
 Curso Superior de Bacharelado em Sistemas de Informação
}} 
\end{center}

\end{tabularx}


% O vfill é um espaço vertical que assume a máxima dimensão possível
% Os vfill's desta página foram utilizados para que o texto ocupe
% toda a folha
\vfill

\LARGE

\textbf{Predição de Crimes Violentos Contra o Patrimônio em Arapiraca e Maceió usando ConvLSTM}

\vfill

\Large

\textbf{Leonardo Victor Fernandes Ferreira} \\
%\textbf{Autor 2}

\vfill

\normalsize
Orientado por:\\
Prof. Dr. Leonardo Melo de Medeiros\\
%Prof. Co-orientaddor


\vfill

\hfill
\parbox{0.5\linewidth}{Trabalho de Conclusão de Curso apresentado como requisito para obtenção do título de Bacharel em Sistemas de Informação.}


\vfill

\large

%\hfill
%\parbox{0.5\linewidth}{\textbf{Trabalho de Conclusão de Curso apresentado ao Curso de Bacharelado em Sistemas de Informação do Instituto Federal de Alagoas como requisito parcial para obtenção do título de bacharel em Sistemas de Informação.}
%
%
%\vfill
%
%\large

%Data
Maceió, AL, agosto de 2016

\end{center}

\end{titlepage}
  

%Espa?amento 1 1/2
\onehalfspacing  

% Paginas introdut?rias (com numeração romana)
\frontmatter

% Lista de conteúdo (sumário, gerado automaticamente)
\addcontentsline{toc}{chapter}{Sumário}   
\tableofcontents  

% Lista de figuras (gerada automaticamente)
\cleardoublepage 
\addcontentsline{toc}{chapter}{Lista de Figuras}  
\listoffigures  

% Lista de tabelas (gerada automaticamente)
\cleardoublepage 
\addcontentsline{toc}{chapter}{Lista de Tabelas}   
\listoftables  

% P?ginas do texto principal (com cabe?alho)
\mainmatter
\pagestyle{headings}

% Para facilitar a organização, foi criado um diretório para cada
% capítulo do documento, pois assim os arquivos das figuras ficam
% classificados por cap?tulos

%%
%% Capitulo 1: Modelo de Capitulo
%%
% Está¡ sendo usando o comando \mychapter, que foi definido no arquivo
% comandos.tex. Este comando \mychapter é essencialmente o mesmo que o
% comando \chapter, com a diferença que acrescenta um \thispagestyle{empty}
% após o \chapter. Isto é necessário para corrigir um erro de LaTeX, que
% coloca um número de página no rodapé de todas as páginas iniciais dos
% capitulos, mesmo quando o estilo de numeração escolhido é outro.
\mychapter{Introdução}
\label{cap:introducao}
% O que escrever
% A introdução deverá ser escrita depois do desenvolvimento do trabalho. Nela deve conter a explicação do que será abordado no trabalho e se necessário um histórico da necessidade do conteudo. Aqui é importante colocar citação direta onde se pode dizer que Segundo \citeasnoun{Bernardete} e ainda escrever seu próprio texto através da leitura de outros trabalhos e fazer uma citação indireta \cite{Bernardete}. Não é muito interessante colocar figuras aqui na introdução embora não seja proibido. Procure escrever de forma resumida, mas faça valer ao leitor o que de importante ele encontrará¡ na leitura, ou seja, aqui deve-se "vender o peixe" do trabalho.
%Sobre esse trabalho
% Descreva o máximo do que se trata o trabalho e sua importáncia.

%NOVO

Os crimes violentos contra o patrimônio (CVP), como roubos, constituem um fenômeno urbano complexo que afeta diretamente a sensação de segurança \cite{camacho22}. Este fenômeno possui uma dinâmica espaço-temporal própria, na qual as ocorrências não se distribuem de forma aleatória, mas tendem a se concentrar de forma persistente em uma pequena parcela do território. Essa regularidade é chamada, na criminologia, de Lei da Concentração do Crime, proposta por David Weisburd (2015), que demonstra que cerca de 50\% dos crimes de uma cidade costumam ocorrer em apenas 4\% a 6\% de seus segmentos de rua.

Nesse sentido, a análise e a predição do crime podem ser tratadas como um problema de previsão espaço-temporal. Onde o espaço urbano é representado por células de área fixa (grids) e o tempo é discretizado em janelas regulares (dias ou semanas) produzindo sequências de mapas de ocorrência. Entretanto, ao adotar a unidade célula–tempo, surge um desafio técnico significativo: a esparsidade de dados. Como a maioria das células não apresenta ocorrências na maior parte do tempo, os registros tornam-se dominados por zeros, caracterizando a predição, na prática, como uma tarefa de previsão de eventos raros sob alto desbalanceamento \cite{wang2019deep}

Diante desse cenário, este trabalho propõe a predição de crimes violentos contra o patrimônio em Arapiraca e Maceió por meio de uma arquitetura de aprendizado profundo do tipo ConvLSTM, adequada para capturar simultaneamente padrões espaciais (no grid) e dependências temporais (na sequência de mapas). Usando como base metodológica, estudo prévio realizado por ... onde foi empregado ConvLSTM para previsão de crime em múltiplas cidades dos Estados Unidos (incluindo Baltimore, Chicago, Los Angeles e Philadelphia) com dados agregados em grids, buscando indicar a viabilidade dessa abordagem para cenários urbanos distintos.


\section{Motivação}
\label{Sec:Motivacao}
%Motivação
Crimes violentos contra o patrimônio impactam diretamente a sensação de segurança, a mobilidade urbana e a atividade econômica, além de pressionarem recursos públicos destinados à prevenção e à resposta  \cite{camacho22,fe2022bad}. Parte fundamental dessa prevenção esta relacionada a coleta e análise de dados relacionados a esses eventos. O Fórum Brasileiro de Segurança Pública divulga anualmente os dados de segurança a nível nacional, tendo o do ano de 2025 apontado transformações importantes nas dinâmicas criminais, com redução de diversos crimes ocorridos no espaço público e, em sentido oposto, aumento de registros ligados a fraudes e golpes (crimes patrimoniais associados ao ambiente virtual). Apesar dessa inversão, os crimes cometidos no espaço físico, especialmente os crimes violentos contra o patrimonio,  continuam a exercer impacto social e econômico, especialmente em centros urbanos como Maceió \cite{fe2022bad,fbsp2025} 

Nesse cenário, é relevante investigar alternativas tecnológicas que apoiem medidas preventivas, principalmente as capazes de estimar o risco futuro por área e por período, oferecendo subsídios para decisões mais precisas para alocação de recursos e direcionamento de ações. Considerando que a ocorrência criminal apresenta dependências no tempo e no espaço e pode se concentrar em hotspots \cite{Weisburd2022}, abordagens de modelagem espaço-temporal tendem a ser mais adequadas do que análises puramente temporais ou apenas descritivas. Assim, avaliar o uso de ConvLSTM com representações em grids, com validação apropriada e discussão de limitações, pode contribuir para a literatura aplicada e para iniciativas locais de planejamento e prevenção. 


\section{Objetivo}
\label{Sec:Objetivo}
%Objetivo
% Descreva o objetivo principal do trabalho. Tente criar um parágrafo resumindo o que seu trabalho fará. Depois seja mais especéfico, pode inclusive criar tópicos para a Sessão \ref{Sec:ObjetivoEspecifico}.

Desenvolver, implementar e avaliar um modelo de aprendizagem  baseado em ConvLSTM para predição de crimes violentos contra o patrimônio, a partir de dados históricos georreferenciados das cidades de Maceió e Arapiraca, organizados em uma representação espacial por grids de modo a capturar padrões de concentração, visando produzir previsões de risco em janelas futuras para apoio analítico a estratégias preventivas.

\subsection{\textbf{Objetivos Específicos}}
\label{Sec:ObjetivoEspecifico}
\begin{itemize}
  \item Organizar e preparar os dados históricos de ocorrências (padronização de datas, georreferenciamento, remoção de inconsistências e tratamento de ausências). ;
  \item Definir a representação em grids, estabelecendo tamanho de célula e granularidade temporal (por exemplo, diário ou semanal), e gerar as sequências de mapas de ocorrência ao longo do período analisado. ;
  \item Caracterizar a esparsidade/desbalanceamento no nível célula–tempo e definir a formulação do alvo (contagem por célula ou variável binária ocorrência/não ocorrência), justificando a escolha. ;
  \item Implementar o modelo ConvLSTM e um pipeline de treino/validação/teste que respeite a ordem temporal (evitando vazamento de informação do futuro). ;
  \item Comparar o ConvLSTM com baselines (ex.: persistência do último mapa, média móvel, ou modelos mais simples), verificando ganhos reais de desempenho. ;
  \item Avaliar o desempenho com métricas adequadas para dados esparsos (incluindo análise por célula e análise agregada por regiões), e discutir limitações, vieses e implicações de uso em contexto real. ;
\end{itemize}

%ANTIGOOO


%Este trabalho tem por objetivo validar se o desempenho e as limitações identificados em estudos prévios sobre previsão e análise espacial do crime se repetem em uma nova realidade urbana (no caso, a cidade de Arapiraca, Alagoas). Em particular, busca-se averiguar se uma das principais limitações observadas em pesquisas internacionais – a baixa precisão preditiva dos modelos – também se manifesta nesse novo contexto. Estudos recentes de previsão criminal em múltiplas cidades reportaram que as ferramentas atuais acertam apenas cerca de 10\% das ocorrências previstas, ou seja, de cada dez alertas de risco de crime, nove acabam não se concretizando. Essa precisão extremamente baixa, constatada de forma consistente nas quatro cidades americanas analisadas por Campedelli et al. (2025), evidencia os limites dos modelos atuais diante da natureza esparsa e volátil dos eventos criminais. Verificar se Arapiraca apresenta desempenho semelhante – isto é, se os algoritmos produzirão muitos falsos positivos e acertos limitados, de modo análogo ao observado nos cenários dos EUA – é fundamental para avaliar a generalização dos achados do estudo original.

%A escolha de Arapiraca como objeto de estudo deve-se à disponibilidade e qualidade dos dados georreferenciados de criminalidade nesse município. Segundo informações da Polícia Militar de Alagoas, houve um esforço de apuração e validação das coordenadas (latitude e longitude) registradas nas ocorrências criminais de Arapiraca, conferindo maior precisão espacial aos dados. Em contrapartida, na capital Maceió tal refinamento não foi realizado com o mesmo rigor, resultando em dados de localização menos confiáveis. Dessa forma, Arapiraca oferece uma base de dados mais consistente para a aplicação da metodologia de análise espacial do crime e para a replicação dos resultados obtidos no estudo original. Em síntese, o objetivo central é testar, no contexto urbano de Arapiraca, se os padrões de desempenho e limitações dos modelos observados em cidades norte-americanas se confirmam, contribuindo assim para avaliar a robustez e alcance geral desses modelos preditivos em uma realidade urbana distinta.

\section{Organização do trabalho}
\label{Sec:Organização}
%Organização do trabalho   
Sempre que precisar referenciar outra parte de seu trabalho use o comando \\ref apontando para o que você colocou no \\label como por exemplo essa Sessão \ref{Sec:Organização}. Aqui você deverá¡ informar como e porque seu trabalho foi organizado, informando o que será abordado em cada capítulo.

% Desenvolvimento
%%
%% Capítulo 2: Desenvolvimento do trabalho
%%

\mychapter{Desenvolvimento}
\label{Cap:desenvolvimento}

Este capítulo descreve a revisão de literatura sobre os fundamentos teóricos utilizados, a metodologia adotada para a coleta, pré-processamento e modelagem dos dados, bem como os detalhes de implementação do pipeline experimental.

\section{Revisão de literatura}
\label{Sec:revisaoliteratura}

%% ================================================================
%% TODO: Seção sobre Criminologia Computacional e Predição de Crime
%% ================================================================
\subsection{TODO -- Predição espacial do crime}
\label{Sec:todo_crime_prediction}

% TODO: Esta subseção precisa ser escrita. Deve cobrir:
%
% 1. Análise espacial do crime e hotspots:
%    - Conceito de hotspot e concentração espacial do crime (já citado na
%      introdução via Weisburd 2015, mas precisa ser aprofundado aqui).
%    - Abordagens tradicionais: Kernel Density Estimation (KDE), modelos
%      de ponto (point-process), ARIMA espacial, etc.
%
% 2. Predição de crime como problema espaço-temporal:
%    - Representação em grids: como o espaço urbano é discretizado em
%      células de área fixa e o tempo em janelas regulares.
%    - Desafio da esparsidade: a maioria das células não apresenta crime
%      na maioria dos períodos, gerando dados altamente desbalanceados.
%    - Formulação do problema: entrada = sequência de mapas passados,
%      saída = mapa de risco futuro.
%
% 3. Aplicações de deep learning na predição de crime:
%    - Evolução dos métodos: de modelos estatísticos clássicos para
%      abordagens baseadas em redes neurais.
%    - Citar e discutir: Wang et al. (2019), Rotaru et al. (2022),
%      Albors Zumel et al. (2025) e outros trabalhos relevantes.
%    - Limitações reportadas na literatura: baixa precisão, vieses,
%      questões éticas (citar Papachristos 2022).
%
% Referências já na .bib que podem ser usadas:
%   \cite{Weisburd2022}, \cite{wang2019deep}, \cite{rotaru2022event},
%   \cite{albors_zumel_deep_2025}, \cite{papachristos_promises_2022}

\textbf{[SEÇÃO PENDENTE]} -- Redigir fundamentação sobre análise espacial do crime, representação espaço-temporal em grids, evolução de métodos de predição criminal (de modelos estatísticos a deep learning) e trabalhos relacionados que empregaram ConvLSTM ou arquiteturas similares para essa tarefa.


\subsection{Redes Neurais Artificiais}

As Redes Neurais Artificiais (RNAs) são sistemas computacionais baseados em modelos matemáticos inspirados no sistema nervoso, buscando aproximar, de forma abstrata, mecanismos de processamento e aprendizagem associados ao cérebro humano \cite{sala2019}. Um marco histórico ocorreu em 1943, quando Warren McCulloch e Walter Pitts propuseram um dos primeiros modelos formais de neurônio artificial através de um modelo matemático que estabeleceu as premissas fundamentais para futuras implementações computacionais \cite{mc1943}. Essas redes operam por processamento paralelo e distribuído, no qual múltiplas unidades simples atuam simultaneamente e de forma interconectada.

Em termos de funcionamento, uma RNA organiza essas unidades em camadas. Cada neurônio recebe valores de entrada, realiza uma combinação ponderada por pesos sinápticos, adiciona um viés e aplica uma função de ativação não linear, gerando um sinal que é propagado para os neurônios da camada seguinte. Essa composição em camadas permite o aprendizado de representações hierárquicas, onde camadas iniciais tendem a capturar padrões mais simples, enquanto camadas posteriores combinam esses padrões para modelar relações mais complexas \cite{gupta2018}.

Além da estrutura em camadas, o desempenho das RNAs depende do processo de treinamento, no qual os pesos e vieses são ajustados para minimizar o erro entre as saídas previstas e os valores reais. Em geral, esse ajuste é realizado por métodos de otimização baseados em gradiente, utilizando o algoritmo de retropropagação do erro (backpropagation) para calcular, de forma eficiente, como cada parâmetro da rede contribui para a função de perda \cite{rumelhart1986learning}. Os diferentes arranjos dessas camadas, bem como a forma como os parâmetros são conectados e compartilhados, dão origem a arquiteturas adequadas a diferentes tipos de dados e tarefas. Redes feedforward processam a informação em um fluxo unidirecional, sendo empregadas em problemas gerais de classificação e regressão. 

Quando os dados apresentam estrutura espacial, as Redes Neurais Convolucionais (CNNs) exploram a organização em grade por meio de operações de convolução e compartilhamento de pesos, tornando-se eficazes para identificar padrões locais e construir representações hierárquicas em imagens e mapas \cite{lecun2015}. Já nos cenários em que a informação é sequencial, as Redes Neurais Recorrentes (RNNs) incorporam um estado interno que é atualizado ao longo do tempo, permitindo a modelagem de dependências temporais entre observações. 
Contudo, RNNs tradicionais podem apresentar limitações para aprender dependências de longo prazo, motivando o desenvolvimento de arquiteturas com mecanismos de memória e portas de controle, como as Long Short-Term Memory (LSTM) \cite{lipton2015}.

%% ================================================================
%% TODO: Seção detalhada sobre LSTM
%% ================================================================
\subsection{TODO -- Long Short-Term Memory (LSTM)}
\label{Sec:todo_lstm}

% TODO: Esta subseção precisa ser escrita. A LSTM é base direta da ConvLSTM
% e merece tratamento próprio. Deve cobrir:
%
% 1. Problema do gradiente desvanecente (vanishing gradient) nas RNNs
%    tradicionais e por que a LSTM foi proposta para resolvê-lo.
%
% 2. Arquitetura da célula LSTM com suas 4 portas:
%    - Porta de esquecimento (forget gate): f_t = sigma(W_f * [h_{t-1}, x_t] + b_f)
%    - Porta de entrada (input gate): i_t = sigma(W_i * [h_{t-1}, x_t] + b_i)
%    - Candidata ao estado celular: g_t = tanh(W_g * [h_{t-1}, x_t] + b_g)
%    - Atualização do estado celular: c_t = f_t * c_{t-1} + i_t * g_t
%    - Porta de saída (output gate): o_t = sigma(W_o * [h_{t-1}, x_t] + b_o)
%    - Estado oculto: h_t = o_t * tanh(c_t)
%
% 3. Incluir diagrama da célula LSTM (recomendado).
%
% 4. Breve menção a variantes (GRU, Peephole LSTM) se desejar.
%
% Referências: \cite{lipton2015}, Hochreiter & Schmidhuber (1997) [adicionar à .bib]

\textbf{[SEÇÃO PENDENTE]} -- Detalhar a arquitetura LSTM: problema do gradiente desvanecente, mecanismo de portas (forget, input, output), equações do estado celular e do estado oculto. Incluir diagrama da célula LSTM.


Essas distinções tornam-se especialmente relevantes em problemas espaço-temporais, nos quais é necessário capturar simultaneamente padrões no espaço e dinâmicas de tempo. Nesse contexto, arquiteturas híbridas como a ConvLSTM combinam convoluções (para modelar as correlações espaciais em grids) com recorrência (para representar sequências de mapas), sendo, portanto, adequadas para tarefas de previsão em dados espaço-temporais \cite{shi2015}.


\subsection{ConvLSTM}

A ConvLSTM (Convolutional Long Short-Term Memory) é uma arquitetura derivada das redes LSTM, proposta por Shi et al. \cite{shi2015}, como já citado, para lidar com sequências espaço-temporais, ou seja, que possuem simultaneamente padrões espaciais e dinâmicas de tempo. Diferentemente da LSTM, em que as transformações entre entrada e estado são realizadas por operações totalmente conectadas, a ConvLSTM substitui essas operações por convoluções, preservando a estrutura espacial dos dados ao longo do tempo. Isso permite que o modelo capture tanto a evolução temporal quanto as correlações espaciais. 

As convoluções consistem na aplicação de pequenos filtros na entrada, também chamados de kernels, que analisam partes pequenas da grade de dados. Esses filtros identificam a concentração, mudanças e semelhanças entre a região analisada e as regiões vizinhas. À medida que percorrem a entrada, as características capturadas são organizadas em diferentes posições do espaço, formando representações que destacam onde determinados padrões ocorrem. 

%% ================================================================
%% TODO: Equações formais da ConvLSTM
%% ================================================================
% TODO: Adicionar as equações formais da ConvLSTM nesta subseção.
% A diferença para a LSTM é que as multiplicações matriciais (W * x)
% são substituídas por convoluções (W ** x):
%
%   i_t = sigma(W_{xi} ** X_t + W_{hi} ** H_{t-1} + b_i)
%   f_t = sigma(W_{xf} ** X_t + W_{hf} ** H_{t-1} + b_f)
%   C_t = f_t . C_{t-1} + i_t . tanh(W_{xc} ** X_t + W_{hc} ** H_{t-1} + b_c)
%   o_t = sigma(W_{xo} ** X_t + W_{ho} ** H_{t-1} + b_o)
%   H_t = o_t . tanh(C_t)
%
% onde ** denota convolução e . denota produto de Hadamard.
% Destacar que X_t, H_t, C_t são tensores 3D (canais x altura x largura),
% o que preserva a estrutura espacial.
% Incluir diagrama comparativo LSTM vs ConvLSTM (recomendado).
%
% Referência: \cite{shi2015}

\textbf{[PENDENTE]} -- Incluir equações formais da ConvLSTM (portas com convolução em vez de multiplicação matricial) e, idealmente, diagrama comparativo LSTM vs ConvLSTM.


%% ================================================================
%% TODO: Técnicas de regularização e treinamento
%% ================================================================
\subsection{TODO -- Técnicas de regularização e otimização}
\label{Sec:todo_regularizacao}

% TODO: Esta subseção precisa ser escrita. Conceitos usados no modelo
% mas que não foram fundamentados na revisão de literatura:
%
% 1. Dropout: desligamento aleatório de neurônios durante o treinamento
%    para prevenir sobreajuste (Srivastava et al., 2014) [adicionar à .bib].
%
% 2. Batch Normalization: normalização das ativações entre camadas
%    para estabilizar e acelerar o treinamento (Ioffe & Szegedy, 2015)
%    [adicionar à .bib].
%
% 3. Funções de ativação: Sigmoid, ReLU, Tanh -- usadas nas portas da
%    ConvLSTM e entre camadas. Breve descrição de cada uma.
%
% 4. Otimizador Adam: combinação de Momentum e RMSProp (Kingma & Ba, 2015)
%    [adicionar à .bib].
%
% 5. Learning rate scheduling: Cosine Annealing (Loshchilov & Hutter, 2017)
%    [adicionar à .bib].

\textbf{[SEÇÃO PENDENTE]} -- Fundamentar Dropout, Batch Normalization, funções de ativação (sigmoid, ReLU, tanh), otimizador Adam e escalonamento de taxa de aprendizado (Cosine Annealing). Todos são usados no modelo mas não estão no referencial.


%% ================================================================
%% TODO: Métricas de avaliação
%% ================================================================
\subsection{TODO -- Métricas de avaliação para classificação binária}
\label{Sec:todo_metricas}

% TODO: Esta subseção precisa ser escrita. As métricas são usadas
% extensivamente na avaliação mas nunca definidas formalmente:
%
% 1. Matriz de confusão: TP, FP, TN, FN.
%
% 2. Precision = TP / (TP + FP)
%
% 3. Recall (Sensibilidade) = TP / (TP + FN)
%
% 4. F1-Score = 2 * (Precision * Recall) / (Precision + Recall)
%
% 5. Threshold de classificação: como logits/probabilidades são
%    convertidos em previsão binária via limiar.
%
% 6. Desbalanceamento de classes: definição formal, impacto nas métricas
%    (accuracy paradox), e estratégias de mitigação:
%    - Ponderação da função de perda (pos_weight)
%    - Focal Loss (Lin et al., 2017) [adicionar à .bib]
%    - Oversampling/undersampling
%
% 7. BCEWithLogitsLoss: combinação de sigmoid + binary cross-entropy
%    com estabilidade numérica.

\textbf{[SEÇÃO PENDENTE]} -- Definir formalmente matriz de confusão, precision, recall, F1-score, threshold de classificação, desbalanceamento de classes e estratégias de mitigação (pos\_weight, focal loss). Também definir BCEWithLogitsLoss.


%% ================================================================
%% TODO: Trabalhos relacionados
%% ================================================================
\subsection{TODO -- Trabalhos relacionados}
\label{Sec:todo_trabalhos_relacionados}

% TODO: Esta subseção precisa ser escrita. Deve discutir os trabalhos
% que servem de base e comparação para este estudo:
%
% 1. Albors Zumel, Tizzoni & Campedelli (2025) -- artigo base deste TCC:
%    - ConvLSTM para crime forecasting em 4 cidades dos EUA
%    - Grid de 0.2 km², previsão 12h à frente, 14 dias de lookback
%    - Comparação com baselines (logistic regression, random forest, LSTM)
%    - Resultados: DL supera baselines, mas precisão geral ainda limitada
%    - Contribuição de mobilidade + sociodemografia
%    \cite{albors_zumel_deep_2025}
%
% 2. Wang et al. (2019) -- deep learning para crime forecasting em tempo real:
%    - Formulação ternária do problema
%    \cite{wang2019deep}
%
% 3. Rotaru et al. (2022) -- predição event-level com viés de enforcement:
%    - Análise de viés policial nos dados e nos modelos
%    \cite{rotaru2022event}
%
% 4. Papachristos (2022) -- promessas e perigos da predição de crime:
%    - Aspectos éticos e limitações
%    \cite{papachristos_promises_2022}
%
% Comparar: dados usados, granularidade, arquitetura, métricas, resultados.

\textbf{[SEÇÃO PENDENTE]} -- Revisar e comparar os trabalhos de referência: Albors Zumel et al. (2025) como base metodológica, Wang et al. (2019) para formulação do problema, Rotaru et al. (2022) para questões de viés, e Papachristos (2022) para aspectos éticos. Comparar dados, granularidade, arquiteturas e resultados obtidos.


%% ================================================================
%%  METODOLOGIA
%% ================================================================
\section{Metodologia}
\label{Sec:metodologia}

Esta seção detalha as etapas do pipeline experimental implementado neste trabalho, desde a obtenção dos dados brutos até a avaliação do modelo treinado. Todas as etapas foram desenvolvidas em Python, utilizando Jupyter Notebooks e bibliotecas como pandas, geopandas, osmnx, shapely, folium, PyTorch e scikit-learn.

\subsection{Dados de origem}
\label{Sec:dados}

Os dados de ocorrências criminais foram obtidos da Polícia Militar de Alagoas por meio de parceria institucional, mediante assinatura de termo de confidencialidade. Em razão dessa restrição, a base de dados original não pode ser disponibilizada publicamente. Os dados correspondem a registros georreferenciados (latitude e longitude) de crimes violentos contra o patrimônio nas cidades de Maceió e Arapiraca, cobrindo o período de 2012 a 2022. O conjunto original contém 137.300 registros com 30 variáveis, incluindo natureza do fato, data/hora, coordenadas geográficas e cidade.

Embora os dados brutos não possam ser divulgados, todo o código-fonte utilizado para o pré-processamento, a construção dos grids, a preparação das sequências e o treinamento do modelo está disponível publicamente em repositório no GitHub\footnote{\url{https://github.com/TODO-inserir-link-do-repositorio}}, permitindo a reprodutibilidade metodológica do trabalho.

\subsection{Pré-processamento dos dados}
\label{Sec:preprocessamento}

O pré-processamento foi realizado em múltiplas etapas encadeadas:

\subsubsection{Carregamento e limpeza inicial}

Os arquivos CSV de entrada foram carregados e concatenados. As colunas de data foram convertidas para o formato datetime (formato \texttt{dd/mm/yyyy HH:MM:SS}) e as coordenadas de latitude e longitude, originalmente com separador decimal em vírgula, foram convertidas para tipo numérico. Registros com data ou coordenadas inválidas (valores ausentes) foram removidos.

\subsubsection{Filtragem por tipo de crime e agrupamento}

Das 23 naturezas de fato distintas presentes nos dados, foram selecionadas 10 categorias relevantes de roubo, que foram agrupadas em 5 macro-categorias:
\begin{itemize}
  \item \textbf{street\_robbery}: roubo a transeunte;
  \item \textbf{vehicle\_robbery}: roubo de veículo (moto, passeio e outros);
  \item \textbf{public\_transport\_robbery}: roubo a transporte coletivo (urbano e rodoviário);
  \item \textbf{commercial\_robbery}: roubo a casa comercial, correspondente bancário e correios;
  \item \textbf{residential\_robbery}: roubo a residência.
\end{itemize}
Após essa filtragem e agrupamento, restaram 119.455 registros, sendo 72.787 em Maceió e 17.421 em Arapiraca, que foram separados em arquivos distintos por cidade.

\subsubsection{Filtragem espacial com polígono municipal}

Para cada cidade, o polígono do limite municipal foi obtido via OpenStreetMap (usando a biblioteca osmnx). Em seguida, realizou-se um spatial join entre os pontos de ocorrência e o polígono, mantendo apenas os registros cujas coordenadas caem efetivamente dentro do município. Para Maceió, 83.512 de 83.612 registros foram mantidos (100 descartados por estarem fora do polígono); para Arapiraca, 19.141 de 19.197 (56 descartados). Os resultados foram visualizados em mapas interativos (folium) para verificação.

\subsection{Construção da representação em grids}
\label{Sec:grids}

\subsubsection{Geração da grade espacial}

Para Maceió, uma grade regular de $39 \times 39$ células foi sobreposta ao bounding box do polígono municipal, resultando em 1.521 células. Cada ocorrência criminal foi atribuída à célula correspondente por meio de verificação de contenção geométrica (point-in-polygon).

% TODO: Calcular e informar a área aproximada de cada célula em km².
% O bounding box de Maceió tem extensão aprox. de ~0.25° lat × ~0.15° lon.
% Com 39 divisões, cada célula tem lado aprox. de X km, resultando em
% uma área de ~Y km². Comparar com o estudo de referência (Albors Zumel
% et al., 2025), que utilizou células de 0,2 km² (0,077 sq. miles).
% Justificar a escolha de 39×39 (ou indicar que foi experimental).
\textbf{[PENDENTE]} -- Informar a área aproximada de cada célula (em km\textsuperscript{2}) e justificar a escolha do tamanho $39 \times 39$. Comparar com o grid de 0,2~km\textsuperscript{2} usado por Albors Zumel et al. (2025).

\subsubsection{Máscara de validade}

Uma máscara binária de $39 \times 39$ foi gerada, indicando quais células do grid possuem interseção com o polígono municipal. Para Maceió, 811 das 1.521 células foram marcadas como válidas. Essa máscara é utilizada na etapa de avaliação para ignorar células que estão fora dos limites da cidade.

\subsubsection{Agregação temporal}

As ocorrências foram contabilizadas por célula e por hora do ano, formando uma matriz esparsa \textit{célula} $\times$ \textit{hora}, para cada tipo de crime e para cada ano. Em seguida, essas matrizes horárias foram agregadas em janelas diárias (granularidade de 24h), resultando em sequências temporais com aproximadamente 365 passos por ano. As contagens dos 5 tipos de crime foram somadas e binarizadas (presença/ausência de crime: $\geq 1 \rightarrow 1$, caso contrário $\rightarrow 0$), produzindo um tensor final de dimensão $4.018 \times 39 \times 39 \times 1$ (dias $\times$ linhas $\times$ colunas $\times$ canais) para Maceió, abrangendo 2012--2022.

\subsection{Preparação para treinamento}
\label{Sec:preparacao}

\subsubsection{Janela deslizante e sub-grids}

Uma janela deslizante de 8 passos (7 dias de lookback + 1 dia de alvo) foi aplicada sobre a sequência temporal. Para cada amostra, em vez de utilizar o grid completo de $39 \times 39$, foram extraídos 5 sub-grids de $16 \times 16$ em posições aleatórias dentro da grade, ampliando o volume de amostras de treinamento e permitindo que o modelo opere sobre regiões locais da cidade.

% TODO: Justificar as seguintes escolhas de hiperparâmetros do pipeline:
% - Por que lookback de 7 dias (e não 14, como no artigo de referência)?
% - Por que sub-grids de 16×16? Qual a relação com o grid total de 39×39?
% - Por que 5 sub-grids por amostra temporal? Isso multiplica os dados por 5.
% - Essas escolhas foram baseadas no artigo de referência, em experimentos
%   preliminares, ou em limitações computacionais?
\textbf{[PENDENTE]} -- Justificar a escolha do lookback de 7 dias, o tamanho do sub-grid ($16 \times 16$) e a quantidade de 5 sub-grids por amostra. Indicar se foram baseados na literatura, em experimentos ou em restrições computacionais.

\subsubsection{Divisão cronológica treino/validação}

A divisão dos dados respeita rigorosamente a ordem temporal para evitar vazamento de informação. Os primeiros 90\% da sequência são utilizados para treinamento e os 10\% finais para validação. Para Maceió, isso resultou em 17.985 amostras de treinamento e 1.925 amostras de validação (cada amostra com formato $7 \times 16 \times 16 \times 1$ para a entrada e $16 \times 16 \times 1$ para o alvo). Os dados foram salvos no formato \texttt{.npz} (arrays comprimidos NumPy). Valores ausentes (NaN) remanescentes foram substituídos por zero antes do treinamento.

% TODO: Discutir as seguintes questões sobre a divisão dos dados:
% - A divisão é apenas treino/validação (90/10). Não há conjunto de teste
%   independente. Isso significa que os resultados reportados são sobre
%   validação, não sobre dados completamente inéditos. Considerar se é
%   necessário adicionar um split de teste (ex.: 80/10/10) ou justificar
%   por que a validação cronológica é suficiente.
% - O split 90/10 coloca ~10 anos em treino e ~1 ano em validação.
%   Indicar quais anos aproximados compõem cada partição.
\textbf{[PENDENTE]} -- Discutir a ausência de conjunto de teste separado (apenas treino/validação). Justificar ou propor um split treino/validação/teste. Indicar quais anos compõem cada partição.


%% ================================================================
%%  MODELO E TREINAMENTO
%% ================================================================
\section{Arquitetura do modelo e treinamento}
\label{Sec:modelo}

\subsection{Arquitetura ConvLSTM implementada}

O modelo foi implementado em PyTorch e consiste em três componentes principais:

\begin{enumerate}
  \item \textbf{ConvLSTM empilhada}: 3 camadas ConvLSTM com 7 canais ocultos cada e kernels de $3 \times 3$. Cada camada ConvLSTM é seguida por ativação ReLU e normalização por lote (BatchNorm3d).
  \item \textbf{Dropout}: camada de dropout com taxa $p = 0{,}8$ aplicada sobre o estado oculto da última camada.
  \item \textbf{Convolução final}: uma convolução $3 \times 3$ com 1 canal de saída (padding ``same''), produzindo o mapa de previsão para o dia seguinte.
\end{enumerate}

O modelo possui 9.262 parâmetros treináveis. A entrada tem formato $(B, 7, 16, 16, 1)$ (batch, sequência, altura, largura, canais) e a saída tem formato $(B, 16, 16)$, representando as probabilidades de ocorrência de crime para cada célula no próximo passo temporal.

% TODO: Incluir justificativa dos hiperparâmetros da arquitetura:
% - Por que 3 camadas ConvLSTM? (profundidade)
% - Por que 7 canais ocultos? (capacidade representacional)
% - Por que kernel 3×3?
% - Por que dropout de 0.8? (valor muito alto, justificar)
% - Essas escolhas foram tuned? Grid search? Baseadas no artigo de referência?
% - Incluir diagrama/figura da arquitetura do modelo (recomendado para TCC).
\textbf{[PENDENTE]} -- Justificar a escolha dos hiperparâmetros da arquitetura (3 camadas, 7 canais, kernel 3$\times$3, dropout 0,8). Indicar se foram baseados no artigo de referência ou em tuning. Incluir diagrama/figura da arquitetura (recomendado).

\subsection{Função de perda e otimização}

Dada a severa esparsidade dos dados ($\sim$2\% das células apresentam crime em qualquer dia), foi utilizada a função de perda \texttt{BCEWithLogitsLoss} com \texttt{pos\_weight} calculado como a razão entre o número de amostras negativas e positivas no conjunto de treinamento ($\approx$48,2 para Maceió). Essa ponderação atribui maior importância aos casos positivos durante o treinamento.

O otimizador utilizado foi o Adam com taxa de aprendizado de $10^{-5}$, e um escalonador do tipo Cosine Annealing com $T_{\max} = 20$ épocas. O treinamento foi conduzido por 200 épocas com batch size de 55, utilizando GPU (NVIDIA GeForce RTX 4070 Laptop GPU) e semente aleatória fixa (42) para reprodutibilidade.

\subsection{Avaliação com tolerância espacial}
\label{Sec:avaliacao}

Além das métricas tradicionais (precision, recall, F1), foi implementada uma avaliação com tolerância espacial, conforme proposto por \citeasnoun{albors_zumel_deep_2025}. Nessa variante, um falso positivo cuja célula vizinha (distância de Chebyshev $\leq 1$) contenha um crime real é reclassificado como acerto espacial (``true positive espacial''). Isso resulta em métricas \textit{precision\_spatial}, \textit{recall\_spatial} e \textit{f1\_spatial}, que reconhecem previsões geograficamente próximas ao evento real.

A avaliação foi realizada sobre o conjunto de validação, varrendo thresholds de classificação de 0,50 a 1,00 (incrementos de 0,01). A máscara de validade municipal foi aplicada para que apenas as células dentro dos limites da cidade contribuíssem para as métricas.





% Resultados e Discussão
%%
%% Capítulo 3: Resultados e Discussão
%%

\mychapter{Resultados e Discussão}
\label{cap:resultados}

Este capítulo apresenta os resultados obtidos com o modelo ConvLSTM treinado para a predição de crimes violentos contra o patrimônio na cidade de Maceió, conforme o pipeline descrito no Capítulo~\ref{Cap:desenvolvimento}. Os experimentos foram conduzidos com semente aleatória fixada em \texttt{seed} $= 42$ para garantir reprodutibilidade, utilizando apenas o canal de ocorrência criminal (modo ``C'') como entrada.

\section{Configuração experimental}
\label{Sec:ConfigExperimental}

O conjunto de dados de Maceió, após a janela deslizante e a extração de sub-grids, resultou em $17.985$ amostras de treinamento e $1.925$ amostras de validação (cada amostra com entrada de formato $7 \times 16 \times 16 \times 1$ e alvo $16 \times 16$). O modelo ConvLSTM implementado possui $9.262$ parâmetros treináveis, arranjados em $3$ camadas empilhadas com $7$ canais ocultos e núcleo convolucional $3 \times 3$, seguidos de uma camada \textit{dropout} com $p = 0{,}8$ e uma convolução final $3 \times 3$ que produz o mapa de probabilidades de ocorrência.

O treinamento foi realizado por $200$ épocas com otimizador Adam (taxa de aprendizado $1 \times 10^{-5}$), \textit{batch size} de $55$ e escalonamento de taxa de aprendizado \texttt{CosineAnnealingLR} com $T_{\max} = 20$. A função de perda adotada foi \texttt{BCEWithLogitsLoss} com \texttt{pos\_weight} $\approx 48{,}18$, calculado como a razão entre células negativas e positivas no conjunto de treinamento, dada a severa esparsidade dos dados ($\sim$2\% das células apresentam crime em qualquer dia).

\section{Curvas de treinamento}
\label{Sec:CurvasTreinamento}

Ao longo das $200$ épocas, observou-se uma redução consistente tanto na perda de treinamento quanto na perda de validação. A perda de treinamento reduziu de $1{,}358$ (época $0$) para $0{,}743$ (época $199$), enquanto a perda de validação reduziu de $1{,}102$ para $0{,}697$ no mesmo intervalo. A ausência de divergência entre as duas curvas, combinada ao fato de a perda de validação se manter sistematicamente abaixo da perda de treinamento, sugere que o modelo não apresentou sobreajuste significativo --- comportamento compatível com o \textit{dropout} de $p = 0{,}8$ aplicado antes da convolução final.

% TODO: inserir figura com a curva de loss (train × val) ao longo das 200 épocas.
\textbf{[PENDENTE]} -- Inserir a figura com a curva de \textit{loss} (treino $\times$ validação) gerada a partir do log de treinamento do notebook \texttt{trainning-maceio.ipynb}.

\section{Métricas por threshold}
\label{Sec:MetricasThreshold}

A saída do modelo, após a aplicação da função sigmoide, fornece uma probabilidade de ocorrência de crime por célula. A conversão dessa probabilidade em classe binária exige a escolha de um \textit{threshold} $\tau$: células com probabilidade $\geq \tau$ são classificadas como positivas. Foram avaliados $51$ valores de \textit{threshold} no intervalo $[0{,}50;\ 1{,}00]$, com passo de $0{,}01$. Para cada \textit{threshold} foram calculadas métricas clássicas de classificação binária (\textit{recall}, \textit{precision} e F1) e suas variantes com tolerância espacial, nas quais um falso positivo é reclassificado como acerto caso exista um crime real em uma célula vizinha (distância de Chebyshev $\leq 1$), conforme descrito na Seção~\ref{Sec:avaliacao}.

\subsection{Métricas tradicionais}
\label{Sec:MetricasTradicionais}

A Tabela~\ref{tab:metricas-tradicionais} apresenta as métricas tradicionais para \textit{thresholds} representativos. O comportamento observado é característico de problemas com classes severamente desbalanceadas: à medida que o \textit{threshold} aumenta, o \textit{recall} cai rapidamente enquanto a \textit{precision} cresce de forma marginal, jamais ultrapassando $\sim$9\%.

\begin{table}[htb]
\centering
\caption{Métricas tradicionais (sem tolerância espacial) por \textit{threshold} para Maceió (\texttt{seed} $= 42$).}
\label{tab:metricas-tradicionais}
\begin{tabular}{cccc}
\hline
\textbf{Threshold} & \textbf{Recall} & \textbf{Precision} & \textbf{F1} \\
\hline
0{,}50 & 0{,}7275 & 0{,}0595 & 0{,}1050 \\
0{,}55 & 0{,}6843 & 0{,}0629 & 0{,}1093 \\
0{,}60 & 0{,}6349 & 0{,}0673 & 0{,}1149 \\
0{,}65 & 0{,}5725 & 0{,}0717 & 0{,}1194 \\
0{,}70 & 0{,}4856 & 0{,}0765 & 0{,}1220 \\
\textbf{0{,}75} & \textbf{0{,}3878} & \textbf{0{,}0838} & \textbf{0{,}1261} \\
0{,}80 & 0{,}2573 & 0{,}0894 & 0{,}1174 \\
0{,}85 & 0{,}1322 & 0{,}0839 & 0{,}0899 \\
0{,}90 & 0{,}0350 & 0{,}0507 & 0{,}0362 \\
\hline
\end{tabular}
\end{table}

O melhor F1 tradicional foi observado em $\tau = 0{,}75$, com valor de $0{,}1261$ (\textit{recall} $= 0{,}3878$ e \textit{precision} $= 0{,}0838$). Mesmo nesse ponto ótimo, a \textit{precision} permanece inferior a $9\%$, indicando que cerca de $91\%$ dos alertas gerados pelo modelo não correspondem a ocorrências reais na célula exata. Esse padrão é consistente com os achados de \citeonline{albors_zumel_deep_2025} em quatro cidades norte-americanas, reforçando a tese de que a métrica célula-a-célula é excessivamente rigorosa para problemas de previsão criminal em grids de alta resolução.

\subsection{Métricas com tolerância espacial}
\label{Sec:MetricasEspaciais}

A Tabela~\ref{tab:metricas-espaciais} apresenta as métricas com tolerância espacial (vizinhança de Chebyshev $\leq 1$) para os mesmos \textit{thresholds}. A mudança de perspectiva --- passar a considerar acerto qualquer previsão que caia em uma das $8$ células vizinhas de um crime real --- produz um salto expressivo em todas as métricas.

\begin{table}[htb]
\centering
\caption{Métricas com tolerância espacial (vizinhança de Chebyshev $\leq 1$) por \textit{threshold} para Maceió (\texttt{seed} $= 42$).}
\label{tab:metricas-espaciais}
\begin{tabular}{cccc}
\hline
\textbf{Threshold} & \textbf{Recall espacial} & \textbf{Precision espacial} & \textbf{F1 espacial} \\
\hline
0{,}50 & 0{,}8637 & 0{,}3313 & 0{,}4540 \\
0{,}55 & 0{,}8443 & 0{,}3464 & 0{,}4642 \\
0{,}60 & 0{,}8199 & 0{,}3595 & 0{,}4719 \\
\textbf{0{,}65} & \textbf{0{,}7771} & \textbf{0{,}3755} & \textbf{0{,}4776} \\
0{,}70 & 0{,}7172 & 0{,}3858 & 0{,}4725 \\
0{,}75 & 0{,}6231 & 0{,}3922 & 0{,}4527 \\
0{,}80 & 0{,}4741 & 0{,}3876 & 0{,}3975 \\
0{,}85 & 0{,}2838 & 0{,}3329 & 0{,}2849 \\
0{,}90 & 0{,}0924 & 0{,}1972 & 0{,}1169 \\
\hline
\end{tabular}
\end{table}

O melhor F1 espacial foi obtido em $\tau = 0{,}65$, atingindo $0{,}4776$ (\textit{recall} espacial $= 0{,}7771$ e \textit{precision} espacial $= 0{,}3755$). Em $\tau = 0{,}50$, o modelo captura $86{,}4\%$ dos crimes reais na vizinhança imediata, ao custo de uma \textit{precision} espacial de $33{,}1\%$. Comparando as Tabelas~\ref{tab:metricas-tradicionais} e \ref{tab:metricas-espaciais}, observa-se que a tolerância espacial multiplica o F1 por um fator de aproximadamente $4\times$ no \textit{threshold} ótimo, evidenciando que o modelo aprende a localizar corretamente vizinhanças de risco, ainda que erre a célula exata.

% TODO: inserir figura com curvas (recall, precision, F1) × threshold para as duas métricas (tradicional e espacial).
\textbf{[PENDENTE]} -- Inserir a figura com as curvas (\textit{recall}, \textit{precision}, F1) $\times$ \textit{threshold} para as duas variantes (tradicional e espacial).

\section{Análise qualitativa}
\label{Sec:AnaliseQualitativa}

% TODO: escolher algumas amostras de validação e apresentar mapas com:
% (a) o alvo real (y_val), (b) a probabilidade predita (antes do threshold),
% (c) a predição binarizada no threshold ótimo (0,65 para F1 espacial).
% As previsões e alvos estão salvos em ./output/maceio/predictions/
% como convlstm_predictions_maceio_42.pkl e convlstm_targets_maceio_42.pkl.
\textbf{[PENDENTE]} -- Apresentar análise qualitativa com mapas lado a lado de: (a) alvo real, (b) probabilidade predita pelo modelo, (c) predição binarizada no \textit{threshold} ótimo ($\tau = 0{,}65$), para amostras representativas do conjunto de validação. As previsões e alvos do último \textit{epoch} encontram-se salvos em \texttt{./output/maceio/predictions/} nos arquivos \texttt{convlstm\_predictions\_maceio\_42.pkl} e \texttt{convlstm\_targets\_maceio\_42.pkl}.

\section{Discussão}
\label{Sec:Discussao}

Os resultados apresentados permitem as seguintes observações:

\begin{itemize}
  \item \textbf{Desbalanceamento dominante.} A esparsidade extrema ($\sim$2\% de células positivas por dia) é o fator estrutural mais relevante. Mesmo com \texttt{pos\_weight} $\approx 48{,}18$, a \textit{precision} tradicional satura abaixo de $9\%$ em todos os \textit{thresholds}, corroborando o diagnóstico de \citeonline{albors_zumel_deep_2025} de que a tarefa célula-a-célula é fundamentalmente mal-posta sob classes tão desbalanceadas.
  \item \textbf{Utilidade prática na vizinhança.} Ao permitir uma tolerância espacial de uma célula, o F1 salta de $\sim$$12\%$ para $\sim$$48\%$. Como cada célula do grid tem $\sim$$0{,}2$~km\textsuperscript{2} de área, uma vizinhança de Chebyshev $\leq 1$ corresponde a uma janela de aproximadamente $3 \times 3$ células ($\sim$$1{,}8$~km\textsuperscript{2}), granularidade coerente com o porte de uma operação de patrulhamento preventivo.
  \item \textbf{Escolha do \textit{threshold}.} Os pontos ótimos das duas métricas são próximos mas distintos ($\tau = 0{,}75$ para F1 tradicional \textit{versus} $\tau = 0{,}65$ para F1 espacial), sugerindo que a calibração deve considerar o tipo de resposta operacional esperada. Valores de $\tau$ no intervalo $[0{,}60;\ 0{,}70]$ oferecem o melhor compromisso entre cobertura e precisão para uso prático.
  \item \textbf{Limitações.} Os valores reportados referem-se a uma única execução com \texttt{seed} $= 42$. Não foram executados \textit{baselines} (persistência, média móvel) para quantificar o ganho efetivo da ConvLSTM, nem foram incorporados canais adicionais (mobilidade, sociodemografia), limitações já elencadas entre os trabalhos futuros.
\end{itemize}

% Conclusões
%%
%% Capítulo 5: Conclusões
%%

\mychapter{Conclusão}
\label{Cap:conclusao}

Este trabalho propôs a aplicação de uma arquitetura ConvLSTM para a predição de crimes violentos contra o patrimônio na cidade de Maceió, utilizando dados georreferenciados de ocorrências entre 2012 e 2022. O pipeline desenvolvido abrangeu desde o pré-processamento dos dados brutos (limpeza, filtragem espacial por polígono municipal, agrupamento de tipos de crime) até a construção de representações espaço-temporais em grids, o treinamento do modelo e a avaliação com métricas tradicionais e com tolerância espacial.

Os resultados preliminares obtidos para Maceió, com semente aleatória 42, indicam que o modelo ConvLSTM é capaz de identificar regiões com maior probabilidade de ocorrência criminal. No threshold de 0,50, o modelo apresentou recall de 72,8\% e precision de 5,9\%, resultando em um F1 de 10,5\%. Esses valores refletem o desafio inerente da esparsidade dos dados: como apenas cerca de 2\% das células apresentam crime em um dado dia, a taxa de falsos positivos é elevada. 

Ao se considerar a avaliação com tolerância espacial (vizinhança de Chebyshev $\leq 1$), os resultados apresentaram melhora expressiva: recall espacial de 86,4\%, precision espacial de 33,1\% e F1 espacial de 45,4\% no threshold de 0,50. O melhor F1 espacial observado foi de 47,8\% no threshold de 0,65, com recall espacial de 77,7\% e precision espacial de 37,5\%. Esses resultados indicam que, embora o modelo não acerte exatamente a célula, uma parcela significativa das previsões incide em vizinhanças imediatas de crimes reais, o que pode ter utilidade prática para o direcionamento de ações preventivas.

Ao longo do treinamento de 200 épocas, observou-se uma redução consistente tanto na perda de treinamento (de 1,358 para 0,743) quanto na perda de validação (de 1,102 para 0,697), sugerindo que o modelo não apresentou sobreajuste significativo apesar do alto dropout aplicado ($p = 0{,}8$).

É importante destacar algumas limitações do presente estudo. A precisão tradicional permaneceu abaixo de 9\% em todos os thresholds, o que significa que a maioria dos alertas gerados pelo modelo não correspondem a ocorrências reais na célula exata. Esse padrão é consistente com os achados da literatura \cite{albors_zumel_deep_2025}, que reporta limitações análogas em cidades norte-americanas. Além disso, o pipeline atual utiliza apenas o canal de crime (modo ``C''), sem incorporar dados sociodemográficos ou de mobilidade, que poderiam contribuir para a melhora do desempenho conforme sugerido por \citeasnoun{albors_zumel_deep_2025}.

O pipeline para Arapiraca encontra-se em fase de pré-processamento: os dados foram filtrados espacialmente (19.141 registros dentro do polígono municipal), mas as etapas de construção de grids, preparação das sequências e treinamento do modelo ainda necessitam ser concluídas.


\section{Trabalhos Futuros}
\label{Sec:TrabalhosFuturos}

Com base nas limitações identificadas e nas possibilidades abertas pelo pipeline desenvolvido, os seguintes trabalhos futuros podem ser realizados:

\begin{itemize}
  \item Completar o pipeline para a cidade de Arapiraca, incluindo a construção de grids, treinamento e avaliação do modelo ConvLSTM, permitindo a comparação entre as duas cidades;
  \item Incorporar canais adicionais de entrada (dados sociodemográficos e de mobilidade) para avaliar se informações contextuais melhoram o desempenho preditivo, conforme o modo ``CMS'' proposto na literatura de referência;
  \item Experimentar diferentes granularidades temporais (12h, 8h) e tamanhos de janela de lookback (14 dias), comparando o impacto na qualidade da predição;
  \item Implementar e comparar modelos baseline (persistência do último mapa, média móvel, regressão logística, Random Forest) para quantificar o ganho efetivo da ConvLSTM;
  \item Avaliar o impacto de diferentes tamanhos de grid e de sub-grid na capacidade preditiva do modelo;
  \item Investigar técnicas de aumento de dados e estratégias alternativas de balanceamento (focal loss, oversampling de amostras positivas) para mitigar o impacto da esparsidade;
  \item Explorar a generalização do modelo por meio de transferência de aprendizado entre cidades (treinamento em Maceió e avaliação em Arapiraca, e vice-versa).
\end{itemize}



\backmatter

% Referências bibliográficas (geradas automaticamente)
%\bibliographystyle{ppgee}
\bibliographystyle{dcu}
\addcontentsline{toc}{chapter}{Referências bibliográficas}
%Arquivo com as bibliografias
\bibliography{bibliografia}


\end{document}